\begin{beginnerstatement}[Kurz erklärt: Gemeinsame Datenregeln]

Eine Schnittstelle ist eine festgelegte Grenze, an der zwei Programmteile
Informationen austauschen. Ein Contract beschreibt dabei, welches Datenformat
beide Seiten erwarten. JSON ist ein textbasiertes Format für strukturierte
Daten, und ein JSON-Schema hält maschinenlesbare Regeln für deren Aufbau fest.
Validierung prüft, ob konkrete Daten diese Regeln einhalten. Frontend und
Backend müssen ihre Formate synchron halten, weil eine einseitige Änderung die
Kommunikation sonst brechen kann. Codegenerierung kann aus einem gemeinsamen
Contract passenden Code für beide Seiten erzeugen und so diese Abstimmung
erleichtern. Im untersuchten Stand existieren Schema und Testbeispiele, sie
werden zur Laufzeit aber nicht automatisch eingebunden; auch eine automatische
Synchronisierung oder Codegenerierung ist noch nicht umgesetzt.

\end{beginnerstatement}
